% =========================================================
% Rational Algebra
% =========================================================

\subsection{Fraction Arithmetic Formulas in $\mathbb{Q}$}
\label{sec:fraction-arithmetic-formulas-q}

After representative-level addition, multiplication, negation, and reciprocals have been defined, the usual fraction formulas can be stated as algebraic rules inside $\mathbb{Q}$.  These results should be read as the familiar symbolic calculus that is justified by the preceding well-definedness theorems.

\begin{definition}[Division and reciprocal in $\mathbb{Q}$]
\label{def:rational-division}
Let $a,b\in \mathbb{Q}$ with $b\neq 0$. We define
\[
\frac{a}{b}
\]
to be the unique rational number $x$ such that
\[
b\cdot x=a.
\]
We also define the reciprocal of $b$ by
\[
b^{-1}:=\frac{1}{b}.
\]
\end{definition}

\begin{remark}
This is not a new primitive operation. It is shorthand for solving the
equation $b\cdot x=a$ inside the field $\mathbb{Q}$.
\end{remark}

\begin{theorem}[Division by One in $\mathbb{Q}$]
\label{thm:rational-division-by-one}
For every $a\in\mathbb{Q}$,
\[
\frac{a}{1}=a.
\]
\end{theorem}

\begin{remark*}[Proof]
\hyperref[prf:rational-division-by-one]{Go to proof.}
\end{remark*}

\begin{theorem}[A Rational Quotient Is Zero Exactly When Its Numerator Is Zero]
\label{thm:rational-quotient-zero-iff-numerator-zero}
For every $a,b\in\mathbb{Q}$ with $b\neq 0$,
\[
\frac{a}{b}=0
\quad\Longleftrightarrow\quad
a=0.
\]
\end{theorem}

\begin{remark*}[Proof]
\hyperref[prf:rational-quotient-zero-iff-numerator-zero]{Go to proof.}
\end{remark*}

\begin{theorem}[Cross-Multiplication Criterion for Rational Equality]
\label{thm:rational-cross-multiplication}
For every $a,b,c,d\in\mathbb{Q}$ with $b\neq 0$ and $d\neq 0$,
\[
\frac{a}{b}=\frac{c}{d}
\quad\Longleftrightarrow\quad
a\cdot d=b\cdot c.
\]
\end{theorem}

\begin{remark*}[Proof]
\hyperref[prf:rational-cross-multiplication]{Go to proof.}
\end{remark*}

\begin{theorem}[Scaling Numerator and Denominator by the Same Nonzero Rational]
\label{thm:rational-fraction-scaling}
For every $a,b,e\in\mathbb{Q}$ with $b\neq 0$ and $e\neq 0$,
\[
\frac{a}{b}=\frac{a\cdot e}{b\cdot e}.
\]
\end{theorem}

\begin{remark*}[Proof]
\hyperref[prf:rational-fraction-scaling]{Go to proof.}
\end{remark*}

\begin{theorem}[Fraction Addition Rule in $\mathbb{Q}$]
\label{thm:rational-fraction-addition}
For every $a,b,c,d\in\mathbb{Q}$ with $b\neq 0$ and $d\neq 0$,
\[
\frac{a}{b}+\frac{c}{d}
=
\frac{a\cdot d+b\cdot c}{b\cdot d}.
\]
\end{theorem}

\begin{remark*}[Proof]
\hyperref[prf:rational-fraction-addition]{Go to proof.}
\end{remark*}

\begin{theorem}[Fraction Subtraction Rule in $\mathbb{Q}$]
\label{thm:rational-fraction-subtraction}
For every $a,b,c,d\in\mathbb{Q}$ with $b\neq 0$ and $d\neq 0$,
\[
\frac{a}{b}-\frac{c}{d}
=
\frac{a\cdot d-b\cdot c}{b\cdot d}.
\]
\end{theorem}

\begin{remark*}[Proof]
\hyperref[prf:rational-fraction-subtraction]{Go to proof.}
\end{remark*}

\begin{theorem}[Fraction Multiplication Rule in $\mathbb{Q}$]
\label{thm:rational-fraction-multiplication}
For every $a,b,c,d\in\mathbb{Q}$ with $b\neq 0$ and $d\neq 0$,
\[
\frac{a}{b}\cdot \frac{c}{d}
=
\frac{a\cdot c}{b\cdot d}.
\]
\end{theorem}

\begin{remark*}[Proof]
\hyperref[prf:rational-fraction-multiplication]{Go to proof.}
\end{remark*}

\begin{theorem}[Negation Rule for Rational Fractions]
\label{thm:rational-fraction-negation}
For every $a,b\in\mathbb{Q}$ with $b\neq 0$,
\[
-\left(\frac{a}{b}\right)
=
\frac{-a}{b}
=
\frac{a}{-b}.
\]
\end{theorem}

\begin{remark*}[Proof]
\hyperref[prf:rational-fraction-negation]{Go to proof.}
\end{remark*}

\begin{theorem}[Reciprocal of a Nonzero Rational Fraction]
\label{thm:rational-fraction-reciprocal}
For every $a,b\in\mathbb{Q}$ with $a\neq 0$ and $b\neq 0$,
\[
\left(\frac{a}{b}\right)^{-1}=\frac{b}{a}.
\]
\end{theorem}

\begin{remark*}[Proof]
\hyperref[prf:rational-fraction-reciprocal]{Go to proof.}
\end{remark*}

\begin{remark}[What this package gives]
These identities are the algebraic rules that make fraction notation usable.
They turn the abstract field structure of $\mathbb{Q}$ into the familiar
symbolic calculus of rational expressions.
\end{remark}
